% Part: first-order-logic
% Chapter: syntax-and-semantics
% Section: main-operator

\documentclass[../../../include/open-logic-section]{subfiles}

\begin{document}

\olfileid{fol}{syn}{mai}

\olsection{સૂત્રનો \printtoken{S}{main operator}}

\begin{explain}
!!a{formula}~$!A$ રચવામાં છેલ્લે વપરાયેલા કારક વિશે વાત કરવી ઘણી
વાર ઉપયોગી બને છે. આ કારકને $!A$નો \emph{મુખ્ય કારક} કહેવાય છે.
સહજ રીતે કહીએ તો તે $!A$નો ``સૌથી બહારનો'' કારક છે. દાખલા તરીકે,
$\lnot !A$નો મુખ્ય કારક~$\lnot$, $(!A\lor !B)$નો મુખ્ય કારક~$\lor$,
વગેરે છે.
\end{explain}


\begin{defn}[!!^{main operator}]
\ollabel{def:main-op}
!!a{formula}~$!A$નો \emph{!!{main operator}} નીચે પ્રમાણે
વ્યાખ્યાયિત થાય છે:
\begin{enumerate}
\item \indcase*{!A}{!A}{$\indfrm$નો કોઈ !!{main operator} નથી.}

\tagitem{prvNot}{\indcase{!A}{\lnot !B}{$\indfrm$નો !!{main operator}
  $\lnot$ છે.}}{}

\tagitem{prvAnd}{\indcase{!A}{(!B \land !C)}{$\indfrm$નો
  !!{main operator} $\land$ છે.}}{}

\tagitem{prvOr}{\indcase{!A}{(!B \lor !C)}{$\indfrm$નો
  !!{main operator} $\lor$ છે.}}{}

\tagitem{prvIf}{\indcase{!A}{(!B \lif !C)}{$\indfrm$નો
  !!{main operator} $\lif$ છે.}}{}

\tagitem{prvIff}{\indcase{!A}{(!B \liff !C)}{$\indfrm$નો
  !!{main operator} $\liff$ છે.}}{}

\tagitem{prvAll}{\indcase{!A}{\lforall[x][!B]}{$\indfrm$નો
  !!{main operator} $\lforall$ છે.}}{}

\tagitem{prvEx}{\indcase{!A}{\lexists[x][!B]}{$\indfrm$નો
  !!{main operator} $\lexists$ છે.}}{}
\end{enumerate}
\end{defn}

દરેક કિસ્સામાં આપણી મનશા સૂત્રમાં દર્શાવેલા !!{main operator}ના ચોક્કસ
\emph{આવર્તન}ની છે. દાખલા તરીકે, સૂત્ર
$((!D\lif !E)\lif(!E\lif !D))$, $(!B\lif !C)$ સ્વરૂપનું છે, જ્યાં
$!B$ એ $(!D\lif !E)$ અને $!C$ એ $(!E\lif !D)$ છે; તેથી $\lif$નું
બીજું આવર્તન !!{main operator} છે.

\begin{explain}
આ બધાં બિન-આણ્વિક !!{formula}sને તેમના !!{main operator}ના આવર્તન
સાથે જોડતા વિધેયની \emph{પુનરાવર્તી} વ્યાખ્યા છે. આપણે
!!{formula}sને જે રીતે અનુમાનાત્મક રીતે વ્યાખ્યાયિત કર્યાં છે તેને
લીધે દરેક !!{formula}~$!A$, \olref{def:main-op}ના કોઈ એક કિસ્સાને
સંતોષે છે. તેથી દરેક બિન-આણ્વિક !!{formula}~$!A$ માટે !!a{main
  operator}નું અસ્તિત્વ સુનિશ્ચિત થાય છે. દરેક !!{formula} આ શરતોમાંથી
માત્ર એકને સંતોષે છે અને દરેક કિસ્સામાં $!A$ જે નાનાં
!!{formula}sમાંથી રચાય છે તે એકમાત્ર રીતે નિર્ધારિત છે; તેથી $!A$ના
!!{main
  operator}નું આવર્તન એકમાત્ર છે અને આમ આપણે વિધેય
વ્યાખ્યાયિત કર્યું છે.
\end{explain}

!!{formula}sનો !!{main operator} જે સંકેત હોય તેના આધારે આપણે તેમને
\olref{tab:main-op}માં આપેલાં નામોથી બોલાવીએ છીએ.
\iftag{defNot,defOr,defAnd,defIf,defIff,defTrue,defFalse,defEx,defAll}
{ જોકે યાદ રાખો કે વ્યાખ્યાયિત કારકો સત્તાવાર રીતે !!{formula}sમાં
આવતા નથી. તેઓ માત્ર સંક્ષિપ્ત રૂપો છે, એટલે સત્તાવાર રીતે તેઓ
સૂત્રના મુખ્ય કારક હોઈ શકે નહિ. તેથી બધાં !!{formula}s અંગેની
સાબિતીઓમાં તેમને અલગથી ધ્યાનમાં લેવાની જરૂર પડતી નથી.}{}
\sourcecorrection{OLSYN-007}{મૂળની \texttt{main-operator.tex}માં
વ્યાખ્યાયિત કારકો અંગેનું સ્મરણપત્ર નિયંત્રિત કરતી
\texttt{\string\iftag} આજ્ઞાનો ત્રીજો દલીલખંડ છૂટી ગયો છે. અહીં ખાલી
ત્રીજો દલીલખંડ પુનઃસ્થાપિત કર્યો છે.}

\begin{table}[!h]
\centering
\begin{tabular}{c | c | c}
!!^{main operator} & !!{formula}નો પ્રકાર & દાખલો\\
\hline
કોઈ નહિ & આણ્વિક (!!{formula}) &
\iftag{prvFalse}{$\lfalse$,}{}
\iftag{prvTrue}{$\ltrue$,}{}
$\Atom{R}{t_1, \dots, t_n}$,
$\eq[t_1][t_2]$\\
$\lnot$ & નિષેધ & $\lnot !A$ \\
$\land$ & સંયોજન & $(!A \land !B)$ \\
$\lor$ & વિયોજન & $(!A \lor !B)$ \\
$\lif$ & !!{conditional} & $(!A \lif !B)$ \\
$\liff$ & !!{biconditional} & $(!A \liff !B)$ \\
$\lforall[][]$ & સાર્વત્રિક (!!{formula})& $\lforall[x][!A]$ \\
$\lexists[][]$ & અસ્તિત્વલક્ષી (!!{formula})& $\lexists[x][!A]$
\end{tabular}
\caption{!!{formula}sના મુખ્ય કારક અને નામો}
\ollabel{tab:main-op}
\end{table}
\sourcecorrection{OLSYN-003}{મૂળની \texttt{main-operator.tex}ની
ત્રણ કોષ્ટક-પંક્તિઓમાં સંયોજન, વિયોજન અને શરતી વિધાનના દાખલાનું
અંતકૌંસ ગણિત પર્યાવરણની બહાર મૂકાયું છે. અહીં ત્રણેય સંપૂર્ણ સૂત્રને
પોતપોતાના એક ગણિત પર્યાવરણમાં પુનઃસ્થાપિત કર્યાં છે.}

\end{document}
